\documentclass[journal]{IEEEtran}

\usepackage{amsmath,amssymb}
\usepackage{graphicx}
\usepackage{booktabs}
\usepackage{cite}

\begin{document}

\title{A Concise IEEE PES Transactions Manuscript Title}

\author{First A. Author,~\IEEEmembership{Member,~IEEE,}
        and Second B. Author,~\IEEEmembership{Senior Member,~IEEE}%
\thanks{Manuscript received Month DD, YYYY; revised Month DD, YYYY.}
\thanks{First A. Author is with the Department of Electrical Engineering,
University Name, City, Country (e-mail: author@example.com).}}

\maketitle

\begin{abstract}
Write a single self-contained paragraph of approximately 150--200 words.
Avoid citations, footnotes, displayed equations, and unexplained abbreviations.
\end{abstract}

\begin{IEEEkeywords}
Active distribution network, demand response, distributed energy resources,
optimal power flow, smart grid.
\end{IEEEkeywords}

\section{Introduction}
State the operational problem, the unresolved technical gap, and the specific
contribution. Tie each claim to evidence that appears later in the paper.

\section{Method}
Define the model, assumptions, notation, and reproducibility-critical settings.

\section{Results}
Refer to each figure and table close to its first discussion, as in Fig.~\ref{fig:example}.

\begin{figure}[!t]
\centering
\fbox{\parbox[c][1.1in][c]{0.92\columnwidth}{\centering
Replace this box with an IEEE-skills PDF or SVG-derived figure.}}
\caption{IEEE one-column figure example. The caption should define plotted
quantities, units, scenarios, and error bars when applicable.}
\label{fig:example}
\end{figure}

\section{Conclusion}
Summarize what was demonstrated, the supporting evidence, and the boundary of
the claim. Do not introduce new results.

\bibliographystyle{IEEEtran}
\bibliography{references}

\end{document}
